\scanpage{49}

séparé $E\otimes F$, \deleted{\unsure{dont|avec|dans}}\inserted{où} le produit scalaire est défini par (3). La norme d'un élément \deleted{\unsure{est notée}} $u\in E\overset{(2)}{\otimes}F$ est notée $\lVert u\rVert_2$\unsure{.} $E\overset{(2)}{\otimes}F$ donc un espace de Hilbert.

\uline{Prop} 1 \unsure{Soient} $E$ et $F$ deux espaces de Hilbert. \deleted{\unsure{1°} \unsure{Si} $x\in E$, $y\in F$, on a $\lVert x\otimes y\rVert_2=\lVert x\rVert\lVert y\rVert$. \unsure{2°}} Soit $(e_i)_{i\in I}$ une base orthonormale dans $E$, $(\unsure{b_j|f_j|b})_{j\in J}$ une base orthonormale dans $F$, alors $(e_i\otimes b_j)_{i,j\in I\times J}$ est une base orthonormale dans $E\overset{(2)}{\otimes}F$.

En effet, d'après (4) le système $(e_i\otimes \unsure{f_j|b_j})$ est orthonormal; \deleted{\unsure{donc pour la totalité} \illegible{2}} \deleted{\unsure{montrons} qu'il est \unsure{dense} dans $E\overset{(2)}{\otimes}F$.} \deleted{\unsure{désignons}} \deleted{En effet, \illegible{2}} \inserted{\unsure{montrons}} \deleted{\unsure{il nous reste à}} que l'espace vectoriel qu'il engendre est dense dans $E\overset{(2)}{\otimes}F$. \deleted{Pour ceci,} Si $E_0$ (resp $F_0$) \deleted{est}\inserted{désigne} l'espace engendré par les $e_i$ (resp. les $f_j$) \deleted{il suffit de} \deleted{\unsure{est évident en effet,}} \deleted{\unsure{fait que|faut que}} $E_0$ est dense dans $E$, $F_0$ dense dans $F$, donc $E_0\otimes F_0$ dense dans $E\otimes F$, donc aussi dans $E\overset{(2)}{\otimes}F$.

Notons \uline{aussi la formule}
\begingroup
\csname tagsleft@true\endcsname
\[
\boxed{\lVert x\otimes y\rVert_2=\lVert x\rVert\lVert y\rVert}
\tag*{(\unsure{7})}
\]
\endgroup
équivalente en effet à \deleted{\unsure{la} \illegible{1}} \deleted{qui résulte \unsure{immédiatement de}} $(x\otimes y,x\otimes y)=(x,x)(y,y)$ \deleted{\unsure{l'égalité qui précède} \illegible{2}}

Notons aussi que l'isomorphisme \inserted{canonique} $\overline{E\otimes F}=\overline{E}\otimes\overline{F}$ (formule (1)) conserve \unsure{manifestement|manifest\textsuperscript{t}} les structures préhilbertiennes, d'où un isomorphisme canonique
\begingroup
\csname tagsleft@true\endcsname
\[
\overline{E\overset{(2)}{\otimes}F}=\overline E\overset{\unsure{(1)|(2)}}{\otimes}\overline F
\tag*{\unsure{(8)|8)}}
\]
\endgroup
% archive sheet number: 50

